% literature.tex  -- written by Harsha
\section{Literature Review}

\subsection{Systolic Computing}
Kung and Leiserson introduced systolic arrays for VLSI matrix
computations~\cite{kung1979systolic}. Kung later characterised them by
rhythmic data flow, high operand reuse, nearest-neighbour communication
and simple replicated cells~\cite{kung1982why}. These are the same
data-parallel principles that underlie SIMD and multiprocessor
design~\cite{culler_pca}.

\subsection{Systolic Arrays in AI Accelerators}
Google's TPU~\cite{jouppi2017tpu} uses a $256\times256$ weight-stationary
array of 8-bit MACs (92~TOPS). It is fed from a unified buffer and a
double-buffered weight FIFO, so that data movement overlaps computation.
Later generations keep the systolic core~\cite{jouppi2021tpuv4i}.
Gemmini~\cite{genc2021gemmini} generates parameterised systolic
accelerators for RISC-V SoCs and supports both OS and WS dataflows.
SCALE-Sim~\cite{samajdar2018scalesim} models array size, dataflow and buffer
effects. All three works show that \emph{data movement and buffering},
not the MAC itself, often decide real throughput. Phase~II of this project
confirms this observation quantitatively.

\subsection{Dataflow Taxonomy}
Sze \emph{et al.}~\cite{sze2017efficient} classify accelerators by which
operand stays in the PE (Table~\ref{tab:dataflow}). Weight-stationary
designs (TPU) preload $B$ and move partial sums. Output-stationary designs
(ShiDianNao~\cite{du2015shidiannao}) keep $C_{ij}$ in place and stream $A$
and $B$. Row-stationary designs (Eyeriss~\cite{chen2017eyeriss}) maximise
reuse in convolutions. We implement OS as the main design and WS as a
measured alternative.

\begin{table}[t]
\centering
\caption{Systolic dataflows}
\label{tab:dataflow}
\setlength{\tabcolsep}{3pt}
\begin{tabular}{@{}lccc@{}}
\toprule
 & \textbf{Weight-st.} & \textbf{Output-st.} & \textbf{Row-st.}\\
\midrule
Resident in PE & $B$ (weights) & $C$ (psum) & filter row \\
Moving operands & $A$, psum & $A$, $B$ & mixed \\
Vertical link width & 32 b (psum) & 8 b ($B$) & mixed \\
Example & TPU~\cite{jouppi2017tpu} & ShiDianNao~\cite{du2015shidiannao} & Eyeriss~\cite{chen2017eyeriss}\\
\bottomrule
\end{tabular}
\end{table}

\subsection{Fixed-Point Arithmetic}
Jacob \emph{et al.}~\cite{jacob2018quantization} showed that 8-bit integer
inference loses little accuracy compared with floating point, while an
8-bit multiplier is far cheaper. Production accelerators therefore pair
8-bit operands with 32-bit accumulators and requantise the
outputs~\cite{jouppi2017tpu}. We follow this scheme with Q4.4 operands, a
Q24.8 accumulator and Q4.4 requantised outputs.

\subsection{Positioning}
Industrial arrays are too large to study in a course project. Our design
is small enough to simulate exhaustively with open-source tools, yet it
exposes the same trade-offs: latency versus throughput, PE count versus
cycles, buffering versus area, clock frequency versus pipeline depth, and
dataflow choice.
